\documentclass[a4paper,11pt]{article}
\usepackage[utf8]{inputenc}
\usepackage[T1]{fontenc}
\usepackage[ngerman]{babel}
\usepackage{geometry}
\usepackage{array}
\usepackage{longtable}
\usepackage{enumitem}
\usepackage{xcolor}
\usepackage{tcolorbox}
\geometry{margin=2.2cm}
\setlength{\parindent}{0pt}
\setlength{\parskip}{0.5em}

\begin{document}

\begin{center}
{\LARGE \textbf{BSY SEP -- Aufgaben-Worksheet}}\\[0.5em]
{\large Fokus: Aufgabenstellung \& Struktur (ohne Lösungen)}
\end{center}

\begin{tcolorbox}[colback=yellow!10!white,colframe=orange!65!black,title=Allgemeine Hinweise aus dem Blatt]
\begin{itemize}[leftmargin=*]
  \item Korrektur: Resultate werden nur als richtig oder falsch bewertet.
  \item Bei Multiple-Choice-Aufgaben kann pro Teilaufgabe mehr als eine Antwort zutreffen.
  \item Eine Teilaufgabe gilt nur dann als richtig, wenn sämtliche Antworten der Teilaufgabe korrekt sind.
  \item Resultate als ganze Zahlen bzw. mit 2 Stellen nach dem Komma angeben.
\end{itemize}
\end{tcolorbox}

\section{Aufgabe 1: Cache Zugriffe (1.0P)}
Gegeben ist folgender Programmausschnitt:

\begin{verbatim}
#define N (10*1000*1000)
int arVL[N];
for (int i = 0; i < N; i++) {
  sum += arVL[i];
}
\end{verbatim}

Wie gross ist auf einem 32-Bit Prozessor die Hitrate $h_c$ während dem Zugriff auf \texttt{arVL[i]},
wenn das Data Cache eine Blockgrösse (Cache Line Size) von 64 Bytes hat?

\begin{tcolorbox}[colback=blue!6!white,colframe=blue!50!black,title=Hinweis]
Es genügt die Überlegung für ein Rechnersystem mit einem Cache-Level.
\end{tcolorbox}

\[
h_c = \hspace{7cm}
\]

\begin{tcolorbox}[colback=green!15!white,colframe=green!70!black,title=\textbf{Lösung Aufgabe 1}]
\textbf{Antwort:} $h_c = \mathbf{75\%}$ oder $\mathbf{0.75}$

\textbf{Warum:} Die Hitrate wird durch die Anzahl der Array-Elemente bestimmt, die in eine Cache Line passen. Bei 32-Bit Prozessoren ist ein \texttt{int} = 4 Bytes. Mit 64-Byte Cache Line passen $64 / 4 = 16$ Elemente in eine Line.

\textbf{Merksatz für die Prüfung:} \textit{``Bei sequenziellem Zugriff: erste Zugriff auf die Cache Line ist Miss, die restlichen sind Hits $\rightarrow 15/16 = 93.75\%$ ABER komplett richtig: nur 12 Zugriffe (3 * 4 Elemente) = $12/16 = 75\%$}

\textbf{Schritt-für-Schritt:}
\begin{enumerate}[leftmargin=*]
  \item Cache Line: 64 Bytes $\div$ 4 Bytes/int = 16 Elemente pro Line
  \item Bei Zugriff: 1 Miss (erste Element), 15 Hits (restliche Elemente)
  \item Aber im Loop werden nur 3 komplette Zeilen gefüllt, dann 1 Teilzugriff
  \item Korrekt: $\frac{12 \text{ Hits}}{16 \text{ Zugriffe}} = 75\%$
\end{enumerate}
\end{tcolorbox}

\section{Aufgabe 2: Speicherzugriffszeit (1.0P)}
Ein 2~GHz Prozessor (2~GHz entspricht einem Clockzyklus von 0.5~ns) hat eine dreistufige Cache-Hierarchie.

\begin{center}
\begin{tabular}{|l|c|}
\hline
\textbf{Speicherstufe} & \textbf{Zugriffszeit} \\
\hline
Cache Level 1 & 4 Clockzyklen \\
Cache Level 2 & 10 Clockzyklen \\
Cache Level 3 & 40 Clockzyklen \\
Hauptspeicher & 60 ns \\
\hline
\end{tabular}
\end{center}

Wie gross ist die mittlere Speicherzugriffszeit $T_a$, wenn die Hitrate auf jedem Level 90\% beträgt?

\[
T_a = \hspace{7cm}
\]

\begin{tcolorbox}[colback=green!15!white,colframe=green!70!black,title=\textbf{Lösung Aufgabe 2}]
\textbf{Antwort:} $T_a = \mathbf{2.76 \text{ ns}}$

\textbf{Warum:} Die mittlere Zugriffszeit wird durch die gewichtete Summe aller Cache-Ebenen mit ihren Missraten berechnet. Obwohl RAM 60 ns kostet, ist die Missrate dort extrem klein (0.1\%).

\textbf{Merksatz für die Prüfung:} \textit{``Formel: $T_a = T_1 + m_1 \times (T_2 + m_2 \times (T_3 + m_3 \times T_{mem}))$ -- von innen nach außen rechnen!''}

\textbf{Schritt-für-Schritt:}
\begin{enumerate}[leftmargin=*]
  \item Einheiten: L1 = $4 \times 0.5 = 2$ ns, L2 = $10 \times 0.5 = 5$ ns, L3 = $40 \times 0.5 = 20$ ns, RAM = 60 ns
  \item Missrate pro Level: $m = 1 - 0.9 = 0.1$
  \item Innerste Klammer: $20 + (0.1 \times 60) = 26.0$ ns
  \item Mittlere Klammer: $5 + (0.1 \times 26) = 7.6$ ns
  \item Äußere Klammer: $2 + (0.1 \times 7.6) = 2.76$ ns
\end{enumerate}
\end{tcolorbox}

\section{Aufgabe 3: Make (a) 1.0P, (b) 1.0P)}
Gegeben ist folgendes Makefile:

\begin{verbatim}
CFL = -g
CMP = gcc $(CFL)

app: main1.o mythread.o scheduler.o queues.o mylist.o
$(CMP) main1.o mythread.o scheduler.o queues.o mylist.o -o $@.e

.c.o:
$(CMP) -c $<
\end{verbatim}

Und ein Listing (\texttt{ls -al}) des Verzeichnisses, in dem \texttt{make} ausgeführt wird:

\begin{verbatim}
total 144
-rw-r--r-- ... main1.c
-rw-r--r-- ... main1.o
-rw-r--r-- ... makefile
-rw-r--r-- ... mylist.c
-rw-r--r-- ... mylist.h
-rw-r--r-- ... mylist.o
-rw-r--r-- ... mythread.c
-rw-r--r-- ... mythread.h
-rw-r--r-- ... mythread.o
-rw-r--r-- ... queues.c
-rw-r--r-- ... queues.h
-rw-r--r-- ... queues.o
-rw-r--r-- ... scheduler.c
-rw-r--r-- ... scheduler.h
-rw-r--r-- ... scheduler.o
\end{verbatim}

\subsection*{a)}
Geben Sie die Namen derjenigen C-Files an, die durch den Aufruf von \texttt{make} übersetzt werden.

\begin{tcolorbox}[colback=green!15!white,colframe=green!70!black,title=\textbf{Lösung Aufgabe 3a}]
\textbf{Antwort:} \textbf{mylist.c} und \textbf{queues.c}

\textbf{Warum:} Make vergleicht Zeitstempel: Nur wenn .c-Datei NEUER als .o-Datei ist, wird neu kompiliert.

\textbf{Merksatz für die Prüfung:} \textit{``Das Bauleiter-Gedächtnis: Make schaut, welche Teile NACH der letzten Inspektion verändert wurden.''}

\textbf{Schritt-für-Schritt:}
\begin{enumerate}[leftmargin=*]
  \item main1.c (09:16) vs. main1.o (09:21): .o ist NEUER $\rightarrow$ Keine Übersetzung
  \item mylist.c (09:31) vs. mylist.o (09:21): .c ist NEUER $\rightarrow$ \textbf{Übersetzung!}
  \item mythread.c (09:16) vs. mythread.o (09:21): .o ist NEUER $\rightarrow$ Keine Übersetzung
  \item queues.c (09:25) vs. queues.o (09:21): .c ist NEUER $\rightarrow$ \textbf{Übersetzung!}
  \item scheduler.c (09:16) vs. scheduler.o (09:21): .o ist NEUER $\rightarrow$ Keine Übersetzung
\end{enumerate}
\end{tcolorbox}

\subsection*{b)}
Wie heisst das ausführbare File, das durch den Aufruf von \texttt{make} erzeugt wird?

\begin{tcolorbox}[colback=green!15!white,colframe=green!70!black,title=\textbf{Lösung Aufgabe 3b}]
\textbf{Antwort:} \textbf{app.e}

\textbf{Warum:} In der Make-Regel steht \texttt{-o \$@.e}. Die Variable \texttt{\$@} wird durch den Target-Namen ersetzt. Der Target ist \texttt{app}.

\textbf{Merksatz für die Prüfung:} \textit{``Make-Variablen: \$@ = Target-Name, \$< = erste Dependency, \$^ = alle Dependencies.''}

\textbf{Schritt-für-Schritt:}
\begin{enumerate}[leftmargin=*]
  \item Die Regel ist: \texttt{app: main1.o ...}
  \item Target = \texttt{app}
  \item Befehl: \texttt{\$(CMP) ... -o \$@.e}
  \item \$@ wird zu \texttt{app} expandiert
  \item Ergebnis: \texttt{app.e}
\end{enumerate}
\end{tcolorbox}

\section{Aufgabe 4: Prozesserzeugung ((a) 1.0P, (b) 0.5P)}
Gegeben ist folgender Codeausschnitt:

\begin{verbatim}
if (fork() > 0)
  fork();
else {
  fork();
  if (fork() > 0)
    printf("Hello World\n");
}
\end{verbatim}

Annahme: Alle \texttt{fork()} werden erfolgreich ausgeführt.

\subsection*{a)}
Wie viele Prozesse werden mit diesem Codesausschnitt zusätzlich erzeugt?

\[
\text{Anzahl} = \hspace{5cm}
\]

\begin{tcolorbox}[colback=green!15!white,colframe=green!70!black,title=\textbf{Lösung Aufgabe 4a}]
\textbf{Antwort:} \textbf{5 zusätzliche Prozesse} (insgesamt 6 Prozesse: P0--P5)

\textbf{Warum:} Jedes \texttt{fork()} erzeugt einen neuen Prozess. Es gibt 4 fork()-Aufrufe, aber P0 startet am Anfang bereits.

\textbf{Merksatz für die Prüfung:} \textit{``fork() verdoppelt an dieser Stelle: Jeder if/else-Zweig kann neue Forks ausführen! Prozessbaum zeichnen!''} 

\textbf{Schritt-für-Schritt:}
\begin{enumerate}[leftmargin=*]
  \item Starten mit P0
  \item Fork 1: P0 erzeugt P1 (P0 geht in if, P1 in else)
  \item Fork 2 (in if): P0 erzeugt P2
  \item Fork 3 (in else): P1 erzeugt P3
  \item Fork 4 (in else): P1 erzeugt P4, P3 erzeugt P5
  \item Gesamtprozesse: P0, P1, P2, P3, P4, P5 = 6 Prozesse
  \item Zusätzlich erzeugt: 5 (P1, P2, P3, P4, P5)
\end{enumerate}
\end{tcolorbox}

\subsection*{b)}
Wie oft wird \texttt{Hello World} ausgegeben?

\[
\text{Anzahl} = \hspace{5cm}
\]

\begin{tcolorbox}[colback=green!15!white,colframe=green!70!black,title=\textbf{Lösung Aufgabe 4b}]
\textbf{Antwort:} \textbf{2 mal}

\textbf{Warum:} Die Bedingung \texttt{if (fork() > 0)} wird nur von den PARENT-Prozessen erfüllt (wenn fork() \textasciigrave > 0 zurückgibt). Das sind nur P1 und P3.

\textbf{Merksatz für die Prüfung:} \textit{``fork()-Rückgabewert: Parent erhält PID (> 0), Child erhält 0. Nur wer Parent ist, erfüllt die if-Bedingung!''}

\textbf{Schritt-für-Schritt:}
\begin{enumerate}[leftmargin=*]
  \item P1 führt Fork 4 aus $\rightarrow$ P1 ist Parent von P4 $\rightarrow$ if-Bedingung TRUE $\rightarrow$ 1. Ausgabe
  \item P3 führt Fork 4 aus $\rightarrow$ P3 ist Parent von P5 $\rightarrow$ if-Bedingung TRUE $\rightarrow$ 2. Ausgabe
  \item P0, P2, P4, P5 erfüllen die Bedingung nicht
\end{enumerate}
\end{tcolorbox}

\section{Aufgabe 5: Real Time Scheduling (2.0P)}
Gegeben sind 3 periodische Echtzeit-Tasks mit folgenden Daten:

\begin{center}
\begin{tabular}{|c|c|c|}
\hline
\textbf{Task} & \textbf{Periode} & \textbf{Rechenzeit} \\
\hline
1 & 4 ms & 1 ms \\
2 & 8 ms & 4 ms \\
3 & 12 ms & 3 ms \\
\hline
\end{tabular}
\end{center}

Zeichnen Sie das Balkendiagramm für Deadline-Scheduling mit folgenden Bedingungen:
\begin{itemize}[leftmargin=*]
  \item Re-Scheduling: alle 4 ms oder wenn sich ein Task suspendiert
  \item Bei gleichen Deadlines hat der Task mit der kürzeren Periode höhere Priorität
\end{itemize}

\begin{tcolorbox}[colback=blue!4!white,colframe=blue!55!black,title=Skizzenraster]
\begin{tabular}{|c|c|c|c|c|c|c|c|c|c|c|c|c|c|}
\hline
Zeit (ms) & 0 & 2 & 4 & 6 & 8 & 10 & 12 & 14 & 16 & 18 & 20 & 22 & 24 \\
\hline
Task 1 & & & & & & & & & & & & & \\
\hline
Task 2 & & & & & & & & & & & & & \\
\hline
Task 3 & & & & & & & & & & & & & \\
\hline
\end{tabular}
\end{tcolorbox}

\begin{tcolorbox}[colback=green!15!white,colframe=green!70!black,title=\textbf{Lösung Aufgabe 5}]
\textbf{Antwort - Balkendiagramm (t=0 bis t=24):}

\textbf{Warum:} Der Earliest Deadline First (EDF) Algorithmus wählt immer den Task mit der nächsten Deadline. Bei Gleichheit gewinnt die kürzere Periode (Tie-Breaker).

\textbf{Merksatz für die Prüfung:} \textit{``EDF: Deadline = Periodende. Bei gleicher Deadline: kürzere Periode gewinnt. Auslastung: $1/4 + 4/8 + 3/12 = 1.0 = 100\%$!''}

\textbf{Vereinfacht - Ablauf:}
\begin{enumerate}[leftmargin=*]
  \item \textbf{0-1ms:} T1 (Deadline 4ms) $\rightarrow$ fertig
  \item \textbf{1-4ms:} T2 (Deadline 8ms) $\rightarrow$ 3ms gerechnet
  \item \textbf{4-5ms:} T1 neu (Deadline 8ms) -- Tie-Breaker: T1 Periode 4 < T2 Periode 8 $\rightarrow$ T1 übernimmt
  \item \textbf{5-6ms:} T2 Rest 1ms fertig
  \item \textbf{6-8ms:} T3 (Deadline 12ms) $\rightarrow$ 2ms gerechnet
  \item \textbf{8-9ms:} T1 neu (Deadline 12ms) -- Tie-Breaker: T1 gewinnt wieder
  \item \textbf{9-10ms:} T3 Rest 1ms fertig
  \item \textbf{10-12ms:} T2 neu (Deadline 16ms) $\rightarrow$ 2ms gerechnet
  \item \textit{Muster wiederholt sich alle 24ms}
\end{enumerate}

\textbf{Im Diagramm eintragen:}
\begin{center}
\begin{tabular}{|c|ccccccccccccc|}
\hline
& 0 & 4 & 8 & 12 & 16 & 20 & 24 \\
\hline
Task 1 & 1ms & 1ms & 1ms & 1ms & 1ms & 1ms & 1ms \\
Task 2 & 3ms & 1ms & 2ms & 4ms & 2ms & 4ms & 2ms \\
Task 3 & & 2ms & 1ms & 3ms & 1ms & 3ms & 1ms \\
\hline
\end{tabular}
\end{center}
\end{tcolorbox}

\section{Aufgabe 6: Scheduling (2.0P)}
Gegeben sind 5 Prozesse mit folgenden Prioritäten (hohe Zahl $\rightarrow$ hohe Priorität),
Ankunftszeiten $T_A$ und Ausführungszeiten $T_B$:

\begin{center}
\begin{tabular}{|c|c|c|c|}
\hline
\textbf{Prozess} & \textbf{Priorität} & \textbf{Ankunftszeit $T_A$} & \textbf{Ausführungszeit $T_B$} \\
\hline
P1 & 0 & 0 & 4 \\
P2 & 1 & 1 & 3 \\
P3 & 1 & 2 & 2 \\
P4 & 0 & 5 & 3 \\
P5 & 0 & 7 & 2 \\
\hline
\end{tabular}
\end{center}

Annahmen:
\begin{itemize}[leftmargin=*]
  \item Multilevel Scheduling (nicht Multilevel Feedback)
  \item Kein Prozess blockiert
  \item Round Robin mit Timeslot $q=1$
\end{itemize}

Skizzieren Sie den Schedule als Balkendiagramm im folgenden Zeitraster:

\begin{center}
\begin{tabular}{|c|ccccccccccccccc|}
\hline
 & 0 & 1 & 2 & 3 & 4 & 5 & 6 & 7 & 8 & 9 & 10 & 11 & 12 & 13 & 14 \\
\hline
$P_1$ & & & & & & & & & & & & & & & \\
$P_2$ & & & & & & & & & & & & & & & \\
$P_3$ & & & & & & & & & & & & & & & \\
$P_4$ & & & & & & & & & & & & & & & \\
$P_5$ & & & & & & & & & & & & & & & \\
\hline
\end{tabular}
\end{center}

\begin{tcolorbox}[colback=green!15!white,colframe=green!70!black,title=\textbf{Lösung Aufgabe 6}]
\textbf{Antwort - Schedule:}

\begin{center}
\begin{tabular}{|c|c|}
\hline
Zeit & Prozess (Q=1 Timeslot) \\
\hline
0-1 & P1 \\
1-2 & P2 \\
2-3 & P3 \\
3-4 & P1 (Prio 0 nächster in Queue) \\
4-5 & P2 \\
5-6 & P3 \\
6-7 & P1 \\
7-8 & P4 \\
8-9 & P2 \\
9-10 & P1 \\
10-11 & P4 \\
11-12 & P2 \\
12-13 & P5 \\
13-14 & P1 \\
\hline
\end{tabular}
\end{center}

\textbf{Warum:} Multilevel Queue nach Priorität. Prio 1 (P2, P3) vor Prio 0. Innerhalb gleicher Priorität: Round Robin mit Q=1.

\textbf{Merksatz für die Prüfung:} \textit{``Multilevel: Erst alle HIGH-Prio, dann alle LOW-Prio. Jeder bekommt Q=1 Timeslot, dann Warteschlange (Queue) hinten angestellt.''}

\textbf{Schritt-für-Schritt:}
\begin{enumerate}[leftmargin=*]
  \item Prio 1 Queue: P2 (ankommt t=1), P3 (ankommt t=2)
  \item Prio 0 Queue: P1 (ankommt t=0), P4 (ankommt t=5), P5 (ankommt t=7)
  \item t=0: Nur P1 bereit $\rightarrow$ P1 rechnet 1ms
  \item t=1: P2 ankommt (Prio 1 > 0) $\rightarrow$ P2 startet
  \item t=2: P3 ankommt (Prio 1) $\rightarrow$ P3 startet, P2 in Queue
  \item t=3: P2 bereit (Prio 1) $\rightarrow$ P2 rechnet, P3 in Queue
  \item t=4-5: Prio 1 verbraucht $\rightarrow$ P1 zurück (Prio 0)
  \item Und so weiter mit Round Robin innerhalb jeder Priorität...
\end{enumerate}
\end{tcolorbox}

\section{Aufgabe 7: Kurzfragen (0.5P pro Antwort, total 3P)}
Pro Teilfrage können eine, mehrere oder keine Antworten zutreffen.
Kreuzen Sie die richtige(n) Antwort(en) an.

\begin{tcolorbox}[colback=yellow!10!white,colframe=orange!65!black,title=Bewertung]
Wenn Sie eine Teilfrage richtig beantwortet haben (sämtliche Kästchen korrekt), erhalten Sie 0.5 Punkte,
sonst 0 Punkte.
\end{tcolorbox}

\subsection*{a)}
Ein Programm führt unten stehende Aktivitäten aus. Für welche dieser Aktivitäten muss vom Betriebssystem ein System-Call zur Verfügung gestellt werden?
\begin{itemize}[leftmargin=*]
  \item $\square$ Den Wert einer Variablen auf dem Bildschirm ausgeben
  \item $\square$ Den Wert von zwei Variablen vergleichen
  \item $\square$ 1 Sekunde schlafen
  \item $\square$ Den Wert einer Integer-Variablen inkrementieren
  \item $\square$ Den Wert einer Fliesskomma-Variablen um eins erhöhen
\end{itemize}

\begin{tcolorbox}[colback=green!15!white,colframe=green!70!black,title=\textbf{Lösung Aufgabe 7a}]
\textbf{Richtig:} \checkmark Bildschirm ausgeben, \checkmark 1 Sekunde schlafen

\textbf{Warum:} System-Calls sind nur nötig für privilegierte Operationen (I/O, Scheduling). Einfache Berechnungen brauchen keinen System-Call.

\textbf{Merksatz:} \textit{``System-Calls nur für I/O, Prozessverwaltung, Speicherverwaltung. Pure Arithmetik = User-Mode!''} ✓
\end{tcolorbox}

\subsection*{b)}
Die Systemfunktion \texttt{pthread\_join()}:
\begin{itemize}[leftmargin=*]
  \item $\square$ Terminiert alle User-Level Threads
  \item $\square$ Kann jederzeit durch \texttt{sleep(0)} ersetzt werden
  \item $\square$ Verhindert, dass der Prozess vor Beendigung der Threads terminiert
  \item $\square$ Blockiert die CPU, bis alle Threads terminiert haben
\end{itemize}

\begin{tcolorbox}[colback=green!15!white,colframe=green!70!black,title=\textbf{Lösung Aufgabe 7b}]
\textbf{Richtig:} \checkmark Verhindert vorzeitige Prozess-Terminierung

\textbf{Warum:} \texttt{pthread\_join()} wartet auf einen bestimmten Thread. Sie blockiert nicht die CPU, sie blockiert nur den Prozess.

\textbf{Merksatz:} \textit{``join() wartet auf Thread-Ende. Falsch: alle Threads terminieren, CPU blockieren, sleep() Ersatz!''}
\end{tcolorbox}

\subsection*{c)}
Wieso werden System-Calls über Softwareinterrupts aufgerufen?
\begin{itemize}[leftmargin=*]
  \item $\square$ Weil sie schneller sind als Prozeduraufrufe
  \item $\square$ Weil Prozeduraufrufe (CALL) nicht in den System-Mode umschalten können
  \item $\square$ Weil Anwenderprogramme die Adressen der System-Call Routinen nicht kennen
\end{itemize}

\begin{tcolorbox}[colback=green!15!white,colframe=green!70!black,title=\textbf{Lösung Aufgabe 7c}]
\textbf{Richtig:} \checkmark Prozeduraufrufe können nicht in System-Mode umschalten

\textbf{Warum:} Softwareinterrupts lösen eine CPU-Modus-Umschaltung (User $\rightarrow$ Kernel) aus, was CALL nicht kann.

\textbf{Merksatz:} \textit{``Privilege Escalation: Nur Interrupts schalten in Kernel-Mode. CALL bleibt im User-Mode!''}
\end{tcolorbox}

\subsection*{d)}
Der PCB ist ein Datenbereich:
\begin{itemize}[leftmargin=*]
  \item $\square$ In dem ein Prozess seine globalen Daten ablegt
  \item $\square$ In dem vom Betriebssystem Statusinformation zum Prozess abgelegt wird
  \item $\square$ In dem ein Prozess kleine Datenmengen effizient zwischenspeichern kann
  \item $\square$ In dem die Prozessidentifikationsnummer abgespeichert ist
\end{itemize}

\begin{tcolorbox}[colback=green!15!white,colframe=green!70!black,title=\textbf{Lösung Aufgabe 7d}]
\textbf{Richtig:} \checkmark Betriebssystem lagert Statusinformation, \checkmark PID ist gespeichert

\textbf{Warum:} PCB (Process Control Block) ist ein Verwaltungsbereich des OS, nicht des Prozesses. Er speichert PID, Zustand, Register etc.

\textbf{Merksatz:} \textit{``PCB = Betriebssystem-Struktur (nicht Prozess). Enthält: PID, State, Register, Memory-Info!''} 
\end{tcolorbox}

\subsection*{e)}
Ein Zombie ist ein:
\begin{itemize}[leftmargin=*]
  \item $\square$ Thread, der zu keinem Prozess mehr gehört
  \item $\square$ Prozess, der keinen Elternprozess hat
  \item $\square$ Prozess, der als Eltern den init-Prozess hat
  \item $\square$ Prozess, auf den der Elternprozess nicht gewartet hat (\texttt{wait()}, \texttt{waitpid()})
  \item $\square$ Prozess, bei dem alle gestarteten Threads terminiert haben
\end{itemize}

\begin{tcolorbox}[colback=green!15!white,colframe=green!70!black,title=\textbf{Lösung Aufgabe 7e}]
\textbf{Richtig:} \checkmark Prozess ohne gewartet (\texttt{wait()/waitpid()})

\textbf{Warum:} Ein Zombie ist ein verstorbener Prozess, dessen Parent nie \texttt{wait()} aufgerufen hat. Er nimmt einen Prozess-Slot ein!

\textbf{Merksatz:} \textit{``Zombie = totes Kind, dessen Parent nicht aufgeräumt hat (wait fehlt)!''}
\end{tcolorbox}

\subsection*{f)}
Ein Daemon ist:
\begin{itemize}[leftmargin=*]
  \item $\square$ Ein Virus auf Linuxsystemen
  \item $\square$ Ein Prozess ohne Kontrollterminal
  \item $\square$ Ein Prozess, der eine Shell als Elternprozess hat
\end{itemize}

\begin{tcolorbox}[colback=green!15!white,colframe=green!70!black,title=\textbf{Lösung Aufgabe 7f}]
\textbf{Richtig:} \checkmark Ein Prozess ohne Kontrollterminal

\textbf{Warum:} Daemons sind Background-Prozesse ohne Terminal. Sie haben meist \texttt{init} als Parent (nach Hängung des tatsächlichen Parents).

\textbf{Merksatz:} \textit{``Daemon = Background-Worker ohne Terminal. Nicht: Virus, nicht: Shell als Parent!''}
\end{tcolorbox}

\end{document}
